\appendix
\section*{Appendix}

\section{Cyclic Proofs}

\subsection{Intuitionistic Martin-L\"of's Inductive Definition System $\LJID$}

We define
an intuitionistic Martin-L\"of's inductive definition system, called $\LJID$.

The language of $\LJID$ is determined by a first-order language with
inductive predicate symbols.
The logical system $\LJID$ is determined by production rules
for inductive predicate symbols.
These production rules mean that 
the inductive predicate denotes 
the least fixed point defined by these production rules.

We assume the first order terms $t,u, \ldots$.
We assume $\forall x$ and $\exists y$ are less tightly connected than other
logical connectives.
To save space,
we sometimes write $Pxy$ and $Fxy$ for $P(x,y)$ and $F(x,y)$.

For example, 
the production rules of the inductive predicate symbol $N$ are
\[
\infer{N0}{}
\qquad
\infer{Ns x}{Nx}
\]
These production rules mean that 
$N$ denotes the smallest set closed under 0 and $s$,
namely the set of natural numbers.

The inference rules of $\LJID$ contains
the introduction rules and the elimination rules for inductive predicates,
determined by the production rules.
These rules describe that the predicate actually denotes the least fixed point.
In particular, the elimination rule describes the induction principle.

For example, the above production rules give the introduction rules
\[
\infer{\Gamma \prove N0}{}
\qquad
\infer{\Gamma \prove Ns x}{\Gamma \prove Nx}
\]
and the elimination rule
\[
\infer{\Gamma,Nt \prove Ft}{
\Gamma \prove F0
&
\Gamma,Fx \prove Fs x
}
\]
This elimination rule describes mathematical induction principle.

\begin{figure*}[t]
{\prooflineskip
\[
\infer[(\Axiom)]{\Gamma,A \prove A}{}
\qquad
\infer[(\Wk)]{\Gamma \prove \Delta}{\Gamma' \prove \Delta'}
(\Gamma' \subseteq \Gamma, \Delta' \subseteq \Delta)
%\\
\qquad
\infer[(\Cut)]{\Gamma \prove \Delta}{
\Gamma \prove F
&
\Gamma,F \prove \Delta
}
%\qquad
\\
\infer[(\Subst)]{\Gamma\theta \prove \Delta\theta}{\Gamma \prove \Delta}
\qquad
\infer[(\neg L)]{\Gamma,\neg F \prove}{\Gamma \prove F}
%\\
\qquad
\infer[(\neg R)]{\Gamma \prove \neg F}{\Gamma,F \prove}
\qquad
\infer[(\lor L)]{\Gamma,F \lor G \prove \Delta}{
\Gamma,F \prove \Delta
&
\Gamma,G \prove \Delta
}
%\qquad
\\
\infer[(\lor R_l)]{\Gamma \prove F \lor G}{
	\Gamma \prove F
}
%\\
\qquad
\infer[(\lor R_r)]{\Gamma \prove F \lor G}{
	\Gamma \prove G
}
\qquad
\infer[(\land L)]{\Gamma,F \land G \prove \Delta}{\Gamma,F,G \prove \Delta}
\qquad
\infer[(\land R)]{\Gamma \prove F \land G}{
\Gamma \prove F
&
\Gamma \prove G
}
\\
\infer[(\imp L)]{\Gamma, F \imp G \prove \Delta}{
	\Gamma \prove F
	&
	\Gamma,G \prove \Delta
}
\qquad
\infer[(\imp R)]{\Gamma \prove F \imp G}{\Gamma,F \prove G}
\qquad
\infer[(\forall L)]{\Gamma,\forall xF \prove \Delta}{\Gamma,F[x:=t] \prove \Delta}
\\
%\]
%\[
\infer[(\forall R)]{\Gamma \prove \forall xF}{\Gamma \prove F}
(x \notin \FV(\Gamma))
\qquad
\infer[(\exists L)]{\Gamma,\exists xF \prove \Delta}{\Gamma,F \prove \Delta}
(x \notin \FV(\Gamma,\Delta))
%\\
\qquad
\infer[(\exists R)]{\Gamma \prove \exists xF}{\Gamma \prove F[x:=t]}
%\qquad
\\
\infer[(=L)]{\Gamma[x:=u],t=u \prove \Delta[x:=u]}{\Gamma[x:=t] \prove \Delta[x:=t]}
\qquad
\infer[(=R)]{\Gamma \prove t=t}{}
\qquad
\infer[(\Ind\ P_j)]{\Gamma,P_j\Vec u \prove F_j\Vec u}{
\hbox{minor premises}
}
\\
\infer[(P\ R)]{\Gamma \prove P\Vec t}{
\Gamma \prove Q_{1}\Vec u_{1} & \ldots & \Gamma \prove Q_{n}\Vec u_{n} 
& \Gamma \prove P_1\Vec t_{1} & \ldots & \Gamma \prove P_{m}\Vec t_{m}
}
\]
}
\vskip -3ex
\caption{Inference Rules}\label{fig:2}
\vskip -1ex
\end{figure*}
The inference rules are given in Figure \ref{fig:2}
where 
for $(P_i\ R)$ we assume the production rule
\[
\infer{P\Vec t}{Q_1\Vec u_1 & \ldots & Q_n\Vec u_n & P_1\Vec t_1 & \ldots & P_m\Vec t_m}
\]
and
for $(\Ind\ P_j)$ we assume a predicate $F_i$ for each $P_i$
and
the minor premises are defined as
\[
\Gamma, Q_{i1}\Vec u_{i1}, \ldots, Q_{in_i}\Vec u_{in_i}, 
F_1\Vec t_{i1}, \ldots, F_{im_i}\Vec t_{im_i}
\prove F_i\Vec t_i
\]
for each production rule
\[
\infer{P_i\Vec t_i}{Q_{i1}\Vec u_{i1} & \ldots & Q_{in_i}\Vec u_{in_i} & P_1\Vec t_{i1} & \ldots & P_{im_i}\Vec t_{im_i}}
\]
Note that the antecedents and the succedents are sets and
the succedent is empty or a formula.

The system $\LJID$ is the same as the system obtained from
classical Martin-L\"of's inductive definition system $\LKID$
defined in \cite{Brotherston11a}
by restricting every sequent to intuitionistic sequents and replacing
$(\imp L)$, $(\lor R)$, and $(\Ind\ P_j)$ accordingly.
The provability of the system $\LJID$ is the same as
that of the natural deduction system given in \cite{Martin-Lof-1971}.

\subsection{Cyclic Proof System $\CLJIDomega$}

An intuitionistic cyclic proof system, called $\CLJIDomega$,
is defined as the system obtained from
classical cyclic proof system $\CLKIDomega$
defined in \cite{Brotherston11a}
by restricting every sequent to intuitionistic sequents and
replacing $(\imp L)$ and $(\lor R)$ in the same way as $\LJID$.
Note that the {\em global trace condition} in $\CLJIDomega$ 
is the same as that in $\CLKIDomega$ (Definition 5.5 of \cite{Brotherston11a}).

Namely, the inference rules of $\CLJIDomega$ are obtained from
$\LJID$ by replacing $(\Ind P_j)$ by
\[
\infer[(\Case P)]{\Gamma, P\Vec u \prove \Delta}{
\hbox{case distinctions}
}
\]
where the case distinctions are
\[
\Gamma,\Vec u=\Vec t,
Q_1\Vec u_1, \ldots, Q_n\Vec u_n, P_1\Vec t_1, \ldots, P_m\Vec t_m
\prove \Delta
\]
for each production rule
\[
\infer{P\Vec t}{Q_1\Vec u_1 & \ldots & Q_n\Vec u_n & P_1\Vec t_1 & \ldots & P_m\Vec t_m}
\]
A cyclic proof in $\CLJIDomega$ is defined by
(1) allowing a bud as an open assumption and
requiring a companion for each bud,
(2) requiring the global trace condition.

The {\em global trace condition} \cite{Brotherston11,Brotherston-phd} is 
the condition that
for every infinite path in the infinite unfolding of a given cyclic proof,
there is a trace that passes
main formulas of case rules
infinitely many times.
The global trace condition ensures that
when we think some measure by counting case rules,
the measure of a bud is smaller than that of the companion.
For example, in the next example the companion (a) uses $Px0y$, but
the bud (a) uses $Px0y$ where $x$ is $x'$ and $x' < x$,
so 
their actual meanings are different
even though they are of the same form.
The global trace condition guarantees the soundness of a cyclic proof system.

An example of a cyclic proof (trivial steps are omitted) is as follows:
\[
\infer[(\Case\ P)]{(a) Px0y \prove x=y}{
\infer{x=0,y=0,Px0y \prove x=y}{
\infer{\prove 0=0}{}
}
&
\infer{Px'0y',x=sx',y=sy' \prove x=y}{
\infer{Px'0y',x=sx',y=sy' \prove sx'=sy'}{
\infer[(\Wk)]{Px'0y',x=sx',y=sy' \prove x'=y'}{
\infer[(\Subst)]{Px'0y' \prove x'=y'}{(a) Px0y \prove x=y}
}}}}
\]
where the mark (a) denotes the bud-companion relation, and
the production rules are
\[
\infer{P0yy}{}
\qquad
\infer{P(sx)y(sz)}{Pxyz}
\]
Note that
the predicate $P$ is addition on natural numbers
and the proof is, essentially, deriving the arithmetic identity $x + 0 = x$.

We call an atomic formula an {\em inductive atomic formula}
when its predicate symbol is an inductive predicate symbol.

\subsection{Cyclic Proofs for Coinductive Predicates}

This subsection shows how we can use a cyclic proof system
to formalize coinductive predicates.
Since a coinductive predicate is a dual of an inductive predicate,
and sequent calculus is symmetric for this dual, we can construct
a cyclic proof system for coinductive predicates.
For example,
for stream predicates
we can define a cyclic proof system $\mu\nu\LK$ from $\CLKIDomega$ as follows:

(1) Add function symbols $\Head$, $\Tail$,
and the pair $\<\ , \ \>$ with
the axioms $\<x,y\>=\<x',y'\> \imp x=x' \land y=y'$ and
$x=\<\Head\ x,\Tail\ x\>$,
and 
a coinductive predicate symbol $P$ with its coproduction rule
\[
\infer[\Co]{P\<y,x\>}{Qyx & Px}
\]
which means $P$ is defined coinductively by this rule.
Note that
$P$ represents
the set of streams $\<x_0, \<x_1, \<x_2, \< \ldots\>\>\>$
such that $Q(x_i,\<x_{i+1}, \<x_{i+2}, \<\ldots\>\>)$ for all $i$.

(2) Add the inference rules $(P\ R)$ and $(\Case\ P)$ in the same way
as $\CLKIDomega$, namely,
\[
\infer[(\Case\ P)]{\Gamma, Pt \prove \Delta}{
\Gamma,t=\<y,x\>, Qyx, Px \prove \Delta
}
\qquad
\infer[(P\ R)]{\Gamma \prove P\<y,x\>,\Delta}{
\Gamma \prove Qyx,\Delta
&
\Gamma \prove Px,\Delta
}
\]

(3) 
We call an atomic formula a {\em coinductive atomic formula}
when its predicate symbol is a coinductive predicate symbol.
We define a {\em cotrace} as a sequence of coinductive atomic formulas
in the succedents of a path such that
two atomic formulas are related by an inference rule in a similar way
to a trace defined in \cite{Brotherston11a}.
The {\em global trace and cotrace condition} is
the condition that
for every infinite path in the infinite unfolding of a given cyclic proof,
the path contains 
either a trace that passes main formulas of case rules infinitely many times,
or a cotrace that passes main formulas of rules $(P\ R)$ infinitely many times.

(4) A cyclic proof is a preproof that satisfies the global trace and cotrace
condition.

Example.
We define the bit stream predicate $\BS$ by the following coproduction rules:
\[
\infer[\Co]{\BS \<y,x\>}{
\Bit\ y
&
\BS\ x
}
\]
where $\Bit$ is an ordinary predicate symbol with the axiom
$\Bit\ x \lequiv x=0 \lor x=1$.
The inference rules from this production rule are:
\[
\infer[(\BS\ R)]{\Gamma \prove \BS\<y,x\>,\Delta}{
\Gamma \prove \Bit\ y,\Delta
&
\Gamma \prove \BS\ x,\Delta
}
\qquad
\infer[(\Case\ \BS)]{\Gamma,\BS\ t \prove \Delta}{
\Gamma, t=\<y,x\>, \Bit\ y, \BS\ x \prove \Delta
}
\]
Then we can show $x=\<0,x\> \prove \BS\ x$, namely,
the zero stream is a bit stream, as follows (trivial steps are omitted):
\[
\infer{(a)\ x=\<0,x\> \prove \BS\ x}{
\infer[(\BS\ R)]{x=\<0,x\> \prove \BS \<0,x\>}{
\infer{x=\<0,x\> \prove \Bit\ 0}{}
&
(a)\ x=\<0,x\> \prove \BS\ x
}}
\]
where (a) denotes the bud-companion relation.

The cyclic proof system $\mu\nu\LK$ is sound for the standard model.
\begin{Th}
If a sequent is provable in $\mu\nu\LK$,
then it is true in the standard model where
a coinductive predicate is interpreted as the greatest fixed point
that satisfies the coproduction rules.
\end{Th}

{\em Proof sketch.}
We add an ordinary predicate symbol $\tilde Q$ with the axiom
$\tilde Qyx \lequiv \neg Qyx$
and add an inductive predicate symbols $\tilde P$ with the production rules
\[
\infer{\tilde Px}{\tilde Q(\Head\ x)(\Tail\ x)}
\qquad
\infer{\tilde Px}{\tilde P(\Tail\ x)}
\]

In the standard model, $P$ is the greatest solution of the equation
\[
Px \lequiv \exists yx'(x=\<y,x'\> \land Qyx' \land Px')
\]
and $\tilde P$ is the least solution of the equation
\[
\tilde Px \lequiv \tilde Q(\Head\ x)(\Tail\ x) \lor \tilde P(\Tail\ x).
\]
By putting $\neg$ on both sides of the equation for $P$
and taking $y$ and $x'$ to be $\Head\ x$ and $\Tail\ x$,
we can show $\neg P$ is a solution of the equation for $\tilde P$.
Hence $\tilde Px \imp \neg Px$ is true.
In the same say by putting $\neg$ on both sides of the equation for $\tilde P$,
and using $x=\<\Head x, \Tail x\>$,
we can show $\neg \tilde P$ is a solution of the equation for $P$.
Hence $\neg \tilde Px \imp Px$ is true.
Therefore $\tilde Px \lequiv \neg Px$ is true.

We define a transformation $(\ )^-$ for a sequent and a proof,
in order to replace $P$ by $\tilde P$.
For a sequent $J$, we define $J^-$ by
replacing $P$ by $\neg \tilde P$ and then
moving an atomic formula $\neg \tilde Pt$ of the antecedent to
$\tilde Pt$ of the succedent and
moving an atomic formula $\neg \tilde Pt$ of the succedent to
$\tilde Pt$ of the antecedent.

Given a cyclic proof $\pi$, we define $\pi^-$ by
replacing each sequent $J$ by $J^-$ and then
replacing $(P\ R)$ by 
\[
\infer[(\Case\ \tilde P)]{\Gamma^-, \tilde P\<y,x'\> \prove \Delta^-}{
\infer{\Gamma^-, \tilde Qyx' \prove \Delta^-}{
\Gamma^- \prove Qyx', \Delta^-
}
&
\Gamma^-, \tilde Px' \prove \Delta^-
}
\ \hbox{(trivial steps are omitted)}
\]
and replacing $(\Case P)$ by 
\[
\infer{\Gamma^- \prove \tilde Px,\Delta^-}{
\infer[(\tilde P\ R)]{\Gamma^-,x=\<y,x'\> \prove \tilde Px,\Delta^-}{
\infer[(\tilde P\ R)]{\Gamma^-,x=\<y,x'\> \prove \tilde Qyx',\tilde Px,\Delta^-}{
\infer{\Gamma^-,x=\<y,x'\> \prove \tilde Qyx',\tilde Px',\Delta^-}{
\Gamma^-,x=\<y,x'\>,Qyx' \prove \tilde Px',\Delta^-
}}}}
\ \hbox{(trivial steps are omitted)}
\]
Then a cotrace in $\pi$ corresponds to a trace in $\pi^-$.
Hence $\pi^-$ is a cyclic proof of $J^-$ in $\CLKIDomega$ when
$\pi$ is a cyclic proof of $J$ in $\mu\nu\LK$.
By the soundness of $\CLKIDomega$, $J^-$ is true in the standard model.
Since $\tilde Px \lequiv \neg Px$ is true,
$J$ is true in the standard model where
$P$ is interpreted as the greatest fixed point.
$\Box$

\section{Equivalence between $\LJID$ and $\CLJIDomega$}

This section studies the equivalence between $\CLJIDomega$ and $\LJID$.

\subsection{Countermodel and Addition of Heyting Arithmetic}

This subsection gives a countermodel and adds arithmetic to the logical systems.

The counterexample given in \cite{Berardi17} also shows that
the equivalence between $\CLJIDomega$ and $\LJID$ does not hold in general,
because the proof of the statement $H$ in \cite{Berardi17} is actually in $\CLJIDomega$,
and $\LJID$ does not prove $H$ since $\LKID$ does not prove $H$.
This gives us the following theorem (it is not new in the sense
\cite{Berardi17} immediately implies it).
\begin{Th}
There are some signature and some set of production rules for which
the provability of $\CLJIDomega$ is not the same as
that of $\LJID$.
\end{Th}

There is a possibility of the equivalence under some conditions.
We will show the equivalence holds by adding arithmetic to both systems.

We add arithmetic to both $\LJID$ and $\CLJIDomega$.
\begin{Def}\rm
$\CLJIDHA$ and $\LJIDHA$
are defined to be obtained from $\CLJIDomega$ and $\LJID$
by adding Heyting arithmetic.
Namely, we add constants and function symbols $0, s, +, \times$,
the inductive predicate symbol $\N$, the productions for $\N$,
and Heyting axioms:
\[
\infer{\N0}{}
\qquad
\infer{\N s x}{\N x}
\qquad
\prove \N x \imp s x \ne 0, \qquad
\prove \N x \land \N y \imp s x=s y \imp x=y, \\
\prove \N x \imp x + 0 = x, \qquad
\prove \N x \land \N y \imp x + s y = s(x+y), \\
\prove \N x \imp x \times 0 = 0, \qquad
\prove \N x \land \N y \imp x \times s y = x \times y + x.
\]
\end{Def}

\subsection{Equivalence Theorem}

In this subsection we state the equivalence theorem and explains proof ideas.

First we assume a new inductive predicate symbol $P'$
for each inductive predicate symbol $P$
and define 
the production rules of $P'$ in the same way as \cite{Berardi17a}.
\begin{Def}\label{def:P'}\rm
We define the production rule of $P'$
\[
\infer{P'\Vec tv}{Q_1\Vec u_1 & \ldots & Q_n\Vec u_n & v > v_1 & P'_1\Vec t_1v_1 & \ldots & v > v_m & P'_m\Vec t_mv_m & \N v}
\]
for each production rule of $P$
\[
\infer{P\Vec t}{Q_1\Vec u_1 & \ldots & Q_n\Vec u_n & P_1\Vec t_1 & \ldots & P_m\Vec t_m}
\]
where $v,v_1,\ldots,v_m$ are fresh variables.
\end{Def}

We write $\LJIDHA+(\Sigma,\Phi)$ for the system
$\LJIDHA$ with the signature $\Sigma$ and the set $\Phi$ of production rules.
Similarly we write $\CLJIDHA+(\Sigma,\Phi)$.
For simplicity, in $\Phi$ we write only $P$ for the set of production rules
for $P$.
We define
$\Sigma_N = \{0,s,+,\times,<,N\}$ and
$ \Phi_N = \{N\}.$
We write $P''$ for $(P')'$.

The next theorem shows
the equivalence of $\LJIDHA$ and $\CLJIDHA$ with signatures.
\begin{Th}[Equivalence of $\LJIDHA$ and $\CLJIDHA$]\label{th:equiv}
Let
$\Sigma = \Sigma_N \cup \{\Vec Q,\Vec P, \Vec P'\}$ and
$\Phi = \Phi_N \cup \{ \Vec P, \Vec P'\}$.
Then
the provability of $\CLJIDHA+(\Sigma,\Phi)$ is the same as
that of $\LJIDHA+(\Sigma,\Phi)$.
\end{Th}

We explain our ideas of proofs of this theorem.
\cite{Berardi17a} shows the equivalence between classical systems
by using classical Podelski-Rybalchenko theorem for induction.
This proof goes well even if we replace classical systems by intuitionistic
systems except that
we have to replace classical Podelski-Rybalchenko theorem for induction
by intuitionistic Podelski-Rybalchenko theorem for induction.
Since we proved intuitionistic Podelski-Rybalchenko theorem for induction
in Theorem \ref{th:Podelski},
by combining them, we can show the equivalence between
$\LJID$ and $\CLJIDomega$.

\section{Proof Transformation}

This section gives the proof of the equivalence.
More detailed discussions are given in \cite{Berardi17b}.
We define proof transformation from $\CLJIDHA$ to $\LJIDHA$.
First we will define stage numbers and path relations,
and then define proof transformation using them.

For notational convenience,
we assume a cyclic proof $\Pi$ in this section.
Let the buds in $\Pi$ be $J_{1i} \ (i \in I)$ and
the companions be $J_{2j} \ (j \in K)$.
Assume $f:I \to K$ such that the companion of a bud $J_{1i}$ is $J_{2,f(i)}$.

\subsection{Stage Numbers for Inductive Definitions}

In this subsection, we define and discuss stage transformation.

We introduce a stage number to each inductive atomic formula 
so that
the argument of the formula comes into the inductive predicate at the stage 
of the stage number.
This stage number will decrease by a progressing trace.
A proof in $\LJIDHA$ will be constructed by 
using the induction on stage numbers.

First we give stage transformation of an inductive atomic formula.
We assume a fresh inductive predicate symbol $P'$ for
each inductive predicate symbol $P$
and we call it a {\em stage-number inductive predicate symbol}.
$P'(\Vec t,v)$ means that
the element $\Vec t$ comes into $P$ at the stage $v$.
We transform $P(\Vec t)$ into $\exists vP'(\Vec t,v)$.
We call a variable $v$ a {\em stage number}
of $\Vec t$ when $P'(\Vec t,v)$.
$P(\Vec t)$ and $\exists vP'(\Vec t,v)$ will become equivalent 
by inference rules introduced by the transformation of production rules.
We call $P'(\Vec t,v)$ a {\em stage-number inductive atomic formula}.

Secondly 
we give stage transformation of a production rule.
We transform the production of $P$
into the production of $P'$ given in Definition \ref{def:P'}.

Next we give the stage transformation of a sequent.
For given fresh variables $\Vec v$,
we transform a sequent $J$ into $J^\circ_{\Vec v}$ defined as follows.
We define $\Gamma^\bullet$ 
as the set obtained from $\Gamma$
by replacing $P(\Vec t)$ by $\exists vP'(\Vec t,v)$.
%
For fresh variables $\Vec v$, 
we define $(\Gamma)^\circ_{\Vec v}$ 
as the sequent obtained from $\Gamma^\bullet$
by replacing the $i$-th element 
of the form $\exists vP'(\Vec t,v)$ in the sequent $\Gamma^\bullet$
by $P'(\Vec t,v_i)$.
%
We define $(\Gamma \prove \Delta)^\bullet$ by $\Gamma^\bullet \prove \Delta^\bullet$,
and define $(\Gamma \prove \Delta)^\circ_{\Vec v}$ by $(\Gamma)^\circ_{\Vec v} \prove \Delta^\bullet$.

We write $(a_i)_{i \in I}$ for the sequence of elements $a_i$ 
where $i$ varies in $I$.
We extend the notion of proofs by allowing open assumptions.
We write $\Gamma \prove_\CLJIDHA \Delta$ with assumptions $(J_i)_{i \in I}$
when there is a proof with assumptions $(J_i)_{i \in I}$ and
the conclusion $\Gamma \prove \Delta$ in $\CLJIDHA$.

\begin{Def}\rm
In a path $\pi$ in a proof,
we define $\Ineq(\pi)$ as the set of the forms
$v > v'$ and $v = v'$ for any stage numbers $v,v'$
eliminated by every case distinction in $\pi$.
\end{Def}

The proof of the next proposition gives 
{\em stage transformation}
of a proof into
a proof of the stage transformation of the conclusion of the original proof.
We write $\Pi^\circ$ for the stage transformation of $\Pi$.

\begin{Prop}[Stage Transformation]\label{prop:birthtrans1}
For any fresh variables $\Vec v$, 
if $\Gamma \prove_\CLJIDHA \Delta$ with assumptions $(\Gamma_i \prove \Delta_i)_{i \in I}$
without any buds,
then
for some fresh variables $(\Vec v_i)_{i \in I}$
we have $(\Gamma)^\circ_{\Vec v} \prove_\CLJIDHA \Delta^\bullet$
with assumptions $(\Ineq(\pi_i),(\Gamma_i)^\circ_{\Vec v_i} \prove \Delta_i^\bullet)_{i \in I}$
without any buds,
where $\pi_i$ is the path from the conclusion to the assumption 
$(\Gamma_i)^\circ_{\Vec v_i} \prove \Delta_i^\bullet$.
\end{Prop}

\subsection{Path Relation}

In this section, we will introduce path relations and discuss them.

We assume a subproof $\Pi_j$ of $\Pi$ 
such that it does not have buds,
its conclusion is $J_{2j}$ and its assumptions are
$J_{1i} \ (i \in I_j)$. 

For $J$ in $\Pi^\circ_j$, we define $\Tilde J$ as $\< v_1,\ldots,v_k \>$ where
$J$ is $\Gamma^\circ_{v_1\ldots v_k} \prove \Delta^\bullet$.

For a path $\pi$ from the conclusion to an assumption in $\Pi_j^\circ$,
we write $\check\pi$ for the corresponding path in $\Pi$.
We extend this notation to a finite composition of $\pi$'s.
By the correspondence $(\check{\ })$,
a stage-number inductive atomic formula in $\Pi^\circ_j$
corresponds to an inductive atomic formula in $\Pi$,
and
a path, a trace, and a progressing trace in $\Pi^\circ_j$
correspond to the same kind of objects in $\Pi$.

\begin{Def}\rm
For a finite composition $\pi$ of paths in $\{ \Pi_j^\circ \ |\ j \in K\}$ such that
$\check\pi$ is a path in the infinite unfolding of $\Pi$,
we define the {\em path relation} $\tilde>_{\pi}$ by
\[
x \tilde>_{\pi} y \equiv
|x|=|\Tilde J_2| \land |y|=|\Tilde J_1| \land
\Land_{F(q_1,q_2)} (x)_{q_2} > (y)_{q_1}
\land
\Land_{G(q_1,q_2)} (x)_{q_2} = (y)_{q_1}
\]
where 
$J_1$ and $J_2$ are the top and bottom sequents of $\pi$ respectively,
$\check J_1$ and $\check J_2$ are those of the path $\check\pi$,
$F(q_1,q_2)$ is that
there is some progressing trace from the $q_2$-th atomic formula in $\check J_2$
to the $q_1$-th atomic formula in $\check J_1$,
$G(q_1,q_2)$ is that
there is some non-progressing trace from the $q_2$-th atomic formula in $\check J_2$ to the $q_1$-th atomic formula in $\check J_1$.
\end{Def}

We define $B_1$ as the set of paths from conclusions to assumptions in 
$\Pi^\circ_j \ (j \in K)$.
We define $B$ as the set of finite compositions of elements in $B_1$ such that
if $\pi \in B$ then $\check\pi$ is a path in the infinite unfolding of $\Pi$.

\begin{Def}\rm\label{def:>}
For $\pi \in B$,
define $x >_{\pi} y$ by
\[
x >_{\pi} y \equiv
(x)_0=j \land (y)_0=f(i) \land  (x)_1 \tilde>_{\pi} (y)_1,
\]
where 
$J_{1i}$ is the top sequent of $\check\pi$, and
$J_{2j}$ is the bottom sequent of $\check\pi$.
\end{Def}

Note that $(\ )_0$ and $(\ )_1$ are operations for a number that represents
a sequence of numbers defined in Section 3.
The first element is a companion number.

\begin{Lemma}\label{lemma:finite}
$\{ >_{\pi} \ |\ \pi \in B \}$ is finite.
\end{Lemma}

{\em Proof.}
Define $C_n$ as $\{ >_{\pi_1\ldots \pi_m} \ |\ m \le n, \pi_i \in B_1 \}$.
Since $>_{\pi}$ is a relation on $N \times N^{\le p}$ where
$p$ is the maximum number of inductive atomic formulas in the antecedents of $\Pi$,
there is $L$ such that $|C_n| \le L$ for all $n$.
Then we have the least $n$ such that $C_{n+1}=C_n$.
Then $|\{ >_{\pi} \ |\ \pi \in B \}|=|C_n|$.
$\Box$

The next lemma is the only lemma that uses the global trace condition.

\begin{Lemma}\label{lemma:globaltrace}
For all $\pi \in B$, there is $n>0$ such that
$\prove_\HA \Ind(U,>_\pi^n)$.
\end{Lemma}

We define $>_\Pi$ as $\bigcup \{ >_{\pi} \ |\ \pi \in B \}$.
Note that $>_\Pi$ is transitive,
since 
the top sequent of $\pi_1$
is the bottom sequent of $\pi_2$ by the first element, and
$((>_{\pi_1}) \circ (>_{\pi_2})) \subseteq (>_{\pi_1\pi_2})$.

\subsection{Proof Transformation}

This section gives proof transformation.

The next lemma shows
we can replace (Case) rules of $\CLJIDHA$ 
by (Ind) rules of $\LJIDHA$.
\begin{Lemma}\label{lemma:case-ind}
If there is a proof
with some assumptions
and without any buds in $\CLJIDHA$,
then
there is a proof of the same conclusions with the same assumptions
in $\LJIDHA$.
\end{Lemma}

The next is a key lemma and shows each bud in a cyclic proof
is provable in $\LJIDHA$, which is proved by using Theorem \ref{th:Podelski}.
\begin{Lemma}\label{lemma:prooftrans}
For every bud $J$ of a proof in $\CLJIDHA$ and fresh variables $\Vec v$,
$(J)^\circ_{\Vec v}$ is provable in $\LJIDHA$.
\end{Lemma}

The next is the main proposition stating that a cyclic proof is transformed
into an $(\LJIDHA)$-proof with stage-number inductive predicates.
\begin{Prop}\label{prop:equiv}
If a sequent $J$ is provable in
$\CLJIDHA+(\Sigma_N \cup \{ \Vec P \}, \Phi_N \cup \{ \Vec P \})$,
then
$J$ is provable in 
$\LJIDHA+(\Sigma_N \cup \{N', \Vec P, \Vec P' \}, \Phi_N \cup \{ \Vec P, \Vec P' \})$
where $N', \Vec P'$ are the stage-number inductive predicates 
of $N, \Vec P$.
\end{Prop}

The next shows conservativity for 
stage-number inductive predicates.
\begin{Prop}[Conservativity of $N'$ and $P''$]\label{prop:conservative}
Let 
$
\Sigma = \Sigma_N \cup \{\Vec Q,\Vec P, \Vec P'\}$,
$\Phi = \Phi_N \cup \{ \Vec P, \Vec P'\}$,
$\Sigma' = \Sigma \cup \{N',\Vec P''\}$, and
$\Phi' = \Phi \cup \{N', \Vec P''\}$.
Then
$\LJIDHA+(\Sigma',\Phi')$
is conservative over
$\LJIDHA+(\Sigma,\Phi)$.
\end{Prop}

\noindent
{\bf Proof of Theorem \ref{th:equiv}.}
(1) $\LJIDHA+(\Sigma,\Phi)$ to $\CLJIDHA+(\Sigma,\Phi)$.

For this claim, we can obtain a proof 
from the proof of Lemma 7.5 in \cite{Brotherston11a}
by restricting every sequent into intuitionistic sequents
and replacing
$\LKID+(\Sigma,\Phi)$ and $\CLKIDomega+(\Sigma,\Phi)$ by
$\LJID+(\Sigma,\Phi)$ and $\CLJIDomega+(\Sigma,\Phi)$ respectively.

(2) $\CLJIDHA+(\Sigma,\Phi)$ to $\LJIDHA+(\Sigma,\Phi)$.

Let
$\Sigma' = \Sigma \cup \{N',\Vec P''\}$ and
$\Phi' = \Phi \cup \{N', \Vec P''\}$.
Assume $J$ is provable in $\CLJIDHA+(\Sigma,\Phi)$.
By Proposition \ref{prop:equiv},
$J$ is provable in $\LJIDHA+(\Sigma',\Phi')$.
By Proposition \ref{prop:conservative},
$J$ is provable in $\LJIDHA+(\Sigma,\Phi)$.
$\Box$

